\cotefile{156-7}
\scanpage{50}
\mspage{49}

\uline{Lemme} 2 Soient $X,Y\in\unsure{M|\mathcal M}$, avec $X\mathrel{>\mkern-13mu<}Y$, $L=\widetilde X\cap\widetilde Y$, $X'\ll X$, $Y'\ll Y$. Posons $X'_L=\widetilde{X'}\cap\mathrm{Omb}(L)$, $Y'_L=\widetilde{Y'}\cap\mathrm{Omb}(L)$. Alors

\[
X'_L\parallel Y'_L\Longleftrightarrow X'\parallel Y'
\]

\deleted{La partie \deleted{\unsure{un peu} délicate est $\Rightarrow$} \inserted{\unsure{principale} dont \unsure{nous avons}} besoin est $\Rightarrow$.} Je n'arrive pas à le prouver en utilisant seulement $M\unsure{\Lambda|\wedge}5$ -- il semble qu'il faut renforcer $M\Lambda5$ par cette forme plus forte, un peu \unsure{draconienne|draconisante|draconiennt} hélas !

Ce lemme est le $M\Lambda4$ de page \unsure{47|41}, Il faut prouver le $M\Lambda5$ de la $=$ page :

\uline{Lemme 3} Soient $X,Y$ avec $X\mathrel{>\mkern-13mu<}Y$, \inserted{$L=\widetilde X\cap\widetilde Y$.} \deleted{\unsure{Si|$\notin$} $L=\widetilde X\cap\widetilde Y$, \unsure{et} $Z\in\Lambda$ \unsure{t.q.|et que} $Z\ll X$,} \deleted{$Z\in\mathrm{omb}_{\Lambda}(X)\cap\mathrm{omb}_{\Lambda}(Y)$} alors $\mathrm{omb}_{\Lambda}(X)\cap\mathrm{omb}_{\Lambda}(Y)=\mathrm{omb}_{\Lambda}(L)$

Je n'y arrive pas non plus. Il faut donc imposer ces deux conditions sur $\parallel_{\Lambda}$ (en termes de la relation $\mathrel{>\mkern-13mu<}_{M}$ qui la définit, cf page précédente).

Donc il faut mettre un $M\Lambda6$ en conséquence. Mais alors, en termes de la proposition p. 42, ça se \unsure{vérifie|vérifiera} tautologiquement : dès qu'on a un système $(M,\leq,\deleted{\unsure{\mathrel{>\mkern-13mu<}}}\ll,\deleted{\unsure{\ll}}\mathrel{>\mkern-13mu<},\Lambda)$, \deleted{\unsure{satisfaisant}} provenant d'un \unsure{atelier|choix} \inserted{spacieux.} \deleted{qui satisfait $\Lambda1,\Lambda2$,} \unsure{Rien|rien} de nouveau ! \unsure{me semble|Il semble} !

\scanpage{51}
\mspage{50}

\uline{Récapitulation}

Les axiomes algébriques essentiels, en termes de la donnée $(\unsure{\mathcal F|\mathscr F|F},\leq,\ll)$, sont récapitulés p. 23. Diagrammes d'implication, concernant les axiomes d'\deleted{\unsure{in|en}}compatibilité-disjonction, voir p. 22 ou p. 45, 46 (avec une \deleted{$\unsure{\Lambda|\wedge}$} partie $\Lambda\subset\unsure{M|\mathcal M}$ dense à la clef, qu'on \unsure{voudra} \unsure{tjrs|tjs} dense -- notion d'atelier $\Lambda$-spacieux ou $\Lambda$-polyédral? \ldots). Sous les noms ``spacieux'' ou ``polyédral'', j'implique maintenant l'axiome $\mathrm{At}\,C(\unsure{\mathrm{spac}|\mathrm{spa}})$ resp. $\mathrm{At}\,\deleted{\unsure{p|po}}C(\mathrm{pol})$.

Dans le cas d'une structure $(M,\leq,\ll,\mathrel{>\mkern-13mu<})$, les données sont \deleted{\unsure{peu|plus}} plus compliquées (et encore, j'ai supprimé $\widetilde{\mathcal F}\inserted{\subset\unsure{P|\mathcal P}(M)}$ comme donnée explicite !), mais le syst. d'axiomes pertinent semble plus simple, \unsure{dans|deux|donc} \deleted{\unsure{certains}} \unsure{puis|pour} certains paquets d'axiomes regroupés.

Il y a lieu de considérer $(M,\leq,\ll,\mathrel{>\mkern-13mu<})$ comme une structure algébrique \unsure{indépendant|indépendante}, on pourrait l'appeler un ``pré-atelier'' (ou ``algèbre \unsure{de|des} multistrates''\deleted{\unsure{)|\}}}, ou ``algèbre de strates''). Elle peut correspondre à \uline{plusieurs} ateliers (réalisés comme des $\unsure{\overline{\mathcal F}|\widetilde{\mathcal F}}\subset P(M)$)

\scanpage{52}
\mspage{51}

parmi lesquels il y en a un le plus grand de tous, formé de \uline{toutes} les parties \inserted{de $\unsure{M|\mathcal M}$} fermées pour $\leq$, et formées d'éléments deux à deux compatibles (pour $\mathrel{>\mkern-13mu<}$).

Je distingue trois descriptions possibles pour les axiomes.

\textcircled{1} \inserted{traduction} Axiomes $\mathrm{At}\,1$ à $\mathrm{At}\,7$ sans plus, première version (cf p. 6\deleted{\unsure{2|a}} -- je laisse tomber les axiomes référant à $\unsure{\mathcal F|\mathscr F|F}$)

\deleted{$\mathrm{At}\,\unsure{M|m}\,1$}

% slanted
\margin{\unsure{ordres|ordre} \unsure{emboîtés|emboîté|multiplicité}}

\[
\text{\uline{Prat} 1}\quad X\leq Y\Rightarrow X\ll Y
\]

% slanted
\margin{transitivité des raff. \unsure{internes}}

\[
\text{\uline{Prat \deleted{\unsure{3|2}}2}}\quad \deleted{\unsure{Y|S}}X\overset{\circ}{\ll}Y\ \text{et}\ Y\overset{\circ}{\ll}Z\ \Rightarrow X\overset{\circ}{\ll}Z
\]

% slanted
\margin{Axiome de la strate \unsure{minimale}}

\[
\text{\uline{Prat \deleted{\unsure{2|3}}3}}\quad
\text{\deleted{\unsure{Si|P}}Pour }X\ll Y,\ \exists\,Y'\leq Y\unsure{,|.}\ \text{avec }X\overset{\circ}{\ll}Y',\ \text{i.e. } X\in\mathrm{Omb}(Y')^{\circ}
\]

\[
\inserted{\text{\uline{Prat 4}}\quad \mathrel{>\mkern-13mu<}\ \text{réflexive, symétrique, et \deleted{\unsure{si}}si }X'\leq X,Y'\leq Y,\quad X\mathrel{>\mkern-13mu<}Y\Rightarrow X'\mathrel{>\mkern-13mu<}Y'}
\]

% slanted
\margin{Unicité de la strate \unsure{minima|minime|minimale}}

\[
\text{\uline{Prat 5}}\quad X\overset{\circ}{\ll}Y,\ X\overset{\circ}{\ll}Y',\quad Y\mathrel{>\mkern-13mu<}Y'\ \Rightarrow Y=Y'
\]

% slanted
\margin{comparaison de $\leq$ et $\ll$ pour strates compatibles}

\[
\text{\uline{Prat 6}}\quad \text{Si }X\mathrel{>\mkern-13mu<}Y,\ \text{alors }X\ll Y\ \Rightarrow X\leq Y
\]

NB seuls les \deleted{deux}\inserted{trois} derniers axiomes font intervenir $\mathrel{>\mkern-13mu<}$

\textcircled{2} Traduction $\mathrm{At}\,1$ - $\mathrm{At}\,7$ sans plus, mais en introduisant de plus une partie $\unsure{\Lambda|\wedge}\subset M$, dont on exigera sans plus qu'elle soit dense (i.e. $\forall X\in\deleted{\unsure{X|F}}M,\ \exists Z\in\Lambda,\ Z\overset{\circ}{\ll}X$, i.e $\mathrm{omb}_{\Lambda}(X)^{\circ}\ne\varnothing$). (Cf page 42)

% slanted
\margin{[NB. Ce syst. d'axiomes est équivalent au \uline{précédent}, quand on fait $\Lambda=M$ ]}

\[
\text{\uline{Prat $\Lambda$ 1}}\quad X\leq Y\ \Rightarrow X\ll Y
\]

\margin{[NB 1, 3, 4 sont comme pour Prat (\unsure{sauf} \unsure{sans|sens|sous} $\Lambda$ pour 3]}

% slanted
\margin{densité de $\Lambda$ (vide si $\Lambda=M$)}

\[
\text{\uline{Prat $\Lambda$ 2}}\quad \forall X\in M,\quad \exists Z\in\Lambda\ \text{avec } Z\overset{\circ}{\ll}X\qquad (\text{i.e. }\mathrm{omb}_{\Lambda}(X)^{\circ}\ne\varnothing)
\]

% slanted
\margin{strate \unsure{minimale|minime} pour les $Z\in\Lambda$}

\[
\text{\uline{Prat $\Lambda$ 3}}\quad \forall X\in M,\quad Z\in\deleted{\unsure{\mathrm{At}\,M|M}}\Lambda,\quad Z\ll X,\ \exists X'\leq X,\ \text{avec }Z\overset{\circ}{\ll}X'
\]

% Prat Lambda 4 is at the foot. Its long arrow points here, under Prat Lambda 3.

\diagram{
\[
\inserted{\begin{aligned}
&\text{\uline{Prat $\Lambda$ 4}}\quad \text{Soient }Y\ll\vcenter{\hbox{\begin{tikzcd}[ampersand replacement=\&,row sep=0pt,cells={nodes={inner sep=0pt}}]
X\\
\unsure{\rotatebox{90}{$\leq$}|\vee\! /}\\
X'
\end{tikzcd}}}\text{ et }Z\in\Lambda\text{ telle que }Z\ll Y,X',\ \text{alors }\exists Y'\ \text{avec}\\
&Z\ll Y',Y'\leq Y,Y'\ll X'.\ \text{Si }Y\leq X,\ \text{\unsure{OPS|ops} }Y'\leq X'.\\
&\bigl[\inserted{\text{i.e. }}\mathrm{omb}_{\Lambda}(Y)\cap\mathrm{omb}_{\Lambda}(X')=\mathrm{omb}_{\Lambda}(Y_{X'})\\
&\text{où }Y_{X'}\overset{\mathrm{déf}}{=}\widetilde Y\cap\mathrm{Omb}(X')\unsure{,|.}
\end{aligned}}
\]
}

\[
\text{\uline{Prat $\Lambda$ 5}}\quad\mathrel{>\mkern-13mu<}\ \text{réflexive et symétrique, \deleted{\unsure{et}}\inserted{si} }X'\leq X,Y'\leq Y,\ \inserted{\text{alors }}X\mathrel{>\mkern-13mu<}Y\Rightarrow X'\mathrel{>\mkern-13mu<}Y'
\]

\[
\begin{aligned}
&\text{\uline{Prat $\Lambda$ 6}}\quad \text{Si } X\mathrel{>\mkern-13mu<}Y,\ \text{et }Z\in\Lambda\ \text{telle que }Z\ll X,Z\ll Y,\ \text{alors}\\
&\exists Z'\in M\ \text{avec } Z'\ll Z\leq X,Y\quad [\ \text{i.e.}\\
&\mathrm{omb}_{\Lambda}(X)\cap\mathrm{omb}_{\Lambda}(Y)=\mathrm{omb}_{\Lambda}(L),\ \text{où }L=\widetilde X\cap\widetilde Y\ ]
\end{aligned}
\]

% slanted
\margin{Axiome des raff. \unsure{induit|induits}}

% slanted
\margin{NB On peut \unsure{bloquer} $\mathrm{Prat}\,\Lambda3,\Lambda4$ en un seul, en exigeant que $\forall Z\ll X$, $\widetilde X^{\unsure{Z|z}}$ a un plus petit élément}

Seuls les deux derniers axiomes font intervenir $\mathrel{>\mkern-13mu<}$.

\scanpage{53}
\mspage{52}

\textcircled{3} Cas des ateliers spacieux ou polyédraux, axiomatique en termes de $(\unsure{M|\mathcal M},\leq,\ll,\parallel_{\unsure{\mathcal L|\mathscr L}})$, où $\parallel_{\mathcal L}$ est une relation ``de disjonction'' sur $\mathcal L={}\text{ens. des él. min. de }(M,\ll)$. (Cf p. 34)

\margin{[Prat dis 1 à Prat dis 5 comme Prat $\unsure{\Lambda|\wedge}$ 1 à Prat $\Lambda$ 5, avec $\Lambda=\mathcal L$]}

% slanted
\margin{dis pour ``disjonction'' ``Prat'' pour ``préatelier''}

\[
\text{\uline{Prat \deleted{\unsure{dis}}\inserted{$\mathcal L$} 1$'$}}\quad X\leq Y\Rightarrow X\ll Y
\]

\[
\text{\uline{Prat \deleted{dis}\inserted{$\mathcal L$} 2$'$}}\quad \text{\deleted{\unsure{Comme}} }\forall X\in M,\quad \mathrm{omb}(X)^{\circ}\ne\varnothing\quad \text{i.e. }\exists x\in\mathcal L,\ x\overset{\circ}{\ll}X
\]

\[
\text{\uline{Prat \deleted{dis}\inserted{$\mathcal L$} 3$'$}}\quad \forall X\in M,\ x\in\mathcal L,\ x\ll X,\quad \exists X'\in M,\quad x\overset{\circ}{\ll}X'\leq X\unsure{,|.}
\]

\uline{Prat \deleted{dis}\inserted{$\mathcal L$} 4$'$} Comme Prat $\deleted{\unsure{A}}\Lambda$ 4 : \unsure{i.e.} Soit

\diagram{
\[
\begin{aligned}
&Y\ll\vcenter{\hbox{\begin{tikzcd}[ampersand replacement=\&,row sep=0pt,cells={nodes={inner sep=0pt}}]
X\\
\unsure{\rotatebox{90}{$\leq$}|\vee\! /}\\
X'
\end{tikzcd}}}\\
&\text{et }x\in\mathcal L,\quad x\ll Y,X',\quad \text{alors }\exists Y'\ \text{avec}\quad Y'\leq Y,\ Y'\ll X',\ x\ll Y'\unsure{,|.}
\end{aligned}
\]
}

\[
\text{i.e.}\qquad \mathrm{omb}(Y)\cap\mathrm{omb}(X')=
\underset{\substack{\parallel\,(\mathrm{déf})\\ \displaystyle\bigcup_{Z\in Y_{X'}}\mathrm{omb}(Z)}}{\mathrm{omb}(Y_{X'})}
\quad \text{où }Y_{X'}=\widetilde Y\cap\mathrm{omb}(X'),
\]

De plus, si $Y\leq X$, \unsure{OPS|ops} $Y'\leq X$, i.e

\[
\mathrm{omb}(Y)\cap\unsure{\mathrm{om}|\mathrm{omb}}(X')=\mathrm{omb}(L)\quad \text{où}\quad L=\widetilde Y\cap\widetilde{X'}.
\]

\uline{Prat \deleted{dis}\inserted{$\mathcal L$} 5$'$} $\parallel_{\mathcal L}$ est \deleted{réflexive}\inserted{symétrique} et antiréflexive (i.e. $x\deleted{=}\parallel y$ implique $y\parallel x$ et $x\ne y$).

\[
\begin{aligned}
&\text{\uline{Prat \deleted{dis}\inserted{$\mathcal L$} 6$'$}}\quad \forall X,Y\in M\ \text{avec }Y\leq X,\quad \deleted{\unsure{xy|x,y}}\ x,y\in\mathcal L,\ y\ll Y,\\
&x\ll X\ \text{et }\deleted{y\in\mathrm{omb}(\unsure{X|Y}}\ x\not\ll Y,\quad (\text{i.e. } y\in\mathrm{omb}(Y),x\in\mathrm{omb}(X)\setminus\mathrm{omb}(Y))\\
&\text{on a }\ x\parallel y.
\end{aligned}
\]

\margin{[Seuls les deux derniers axiomes font intervenir $\parallel$]}

Ce sont les \unsure{$\overline{\mathrm{ms}}$|$\Longleftrightarrow$} axiomes que Prat $\Lambda$, à l'exception du dernier Prat dis 6 (et à l'exception du fait que Prat dis 5 prend une forme plus simple) -- et qui \unsure{remplacent|remplace} \deleted{\unsure{les}} \unsure{douzaines|dizaines|dizaine|derniers} d'axiomes, \unsure{correspondant|correspond.t|en exprimant} \unsure{$\mathrel{>\mkern-13mu<}$|et} en termes d'

% slanted
\margin{Donc le \uline{\unsure{$\overline{\mathrm{m}}$|$\overline{\mathrm{me}}$}} préatelier, permet de construire \unsure{à la fois|à la|à} un atelier polyédral et un atelier spacieux (le plus \unsure{souvent|souvt|svt} différ\unsure{ents|ts}\ldots)}

On définira :

\[
X\mathrel{>\mkern-13mu<}Y\Longleftrightarrow
\left\{\begin{aligned}
&\forall\deleted{\unsure{a}}x\in\mathrm{omb}(X)\setminus\mathrm{omb}(\deleted{\unsure{X}}L),y\in\mathrm{omb}(Y)\setminus\mathrm{omb}(L)\\
&\qquad(\text{où }L=\widetilde X\cap\widetilde Y)\ \text{on a }\ x\parallel y\\
&[+\ L=\varnothing\ \text{ou } L\ \text{de la forme }\widetilde Z,\ \text{dans le}\\
&\hspace{16em}\text{cas ``polyédral''}\ ]
\end{aligned}\right.
\]

\scanpage{54}
\mspage{53}

Mais dans le cas polyédral, il faut un axiome suppl.

\[
\begin{aligned}
&\text{\uline{Prat pol}}\quad \forall X\in\unsure{M|\mathcal M},\quad \widetilde X\ \text{est polyédral,}\quad \text{i.e. si } \deleted{\unsure{X}}Y,Z\in M\ \text{et}\\
&\left\lvert\begin{aligned}
&\text{si }\{X,Y\}\ \text{est majoré et minoré dans }M,\ \text{alors }\mathrm{Inf}^{\leq}(\deleted{\unsure{X}}Y,\deleted{\unsure{Y}}Z)\ \text{existe}\\
&(\text{i.e. }\widetilde X\cap\widetilde Y\ \text{a un plus grand élément}).
\end{aligned}\right.
\end{aligned}
\]

\uline{NB} Alors $\mathrm{Sup}^{\leq}(\deleted{\unsure{X}}Y,\deleted{\unsure{Y}}Z)$ existe aussi \uline{dans} $\widetilde X$ (mais prob. \unsure{bien|sûr} \uline{pas} dans $M$\ldots)

\uline{Remarque} Soit $(M,\leq)$ un ens. ordonné, et \unsure{prenons|posons} $(\ll)\Leftrightarrow(\leq)$, $(x\parallel_{\unsure{\mathcal L|\mathscr L}}y)\Leftrightarrow(x\ne y)$. Alors tous les axiomes sont satisfaits, sauf le seul axiome de divisibilité $\mathrm{Prat}\,\mathcal L\,2'$. Je pense p.ex. au cas où $M={}\text{ensemble des simplexes}$ d'une maquette simpliciale, $(\unsure{\mathscr S|S},M\subset\unsure{P|\mathcal P}_{f}(\mathscr S))$, où les figures sont les sous-maquettes simpliciales, \unsure{correspondants|correspondant} aux sous-ensembles fermés de $(M,\leq(\overset{\mathrm{déf}}{=}\subset))$.

En termes de l'axiomatique \textcircled{1} ou \textcircled{2}, on prendrait $X\ll Y\deleted{\unsure{=} }\Longleftrightarrow X\leq Y$, et $X\mathrel{>\mkern-13mu<}Y$ tjrs satisfaits. \deleted{Alors $X\overset{\circ}{\ll}Y$ signifie $X=Y$.}

Appelons ``trivial'' un \inserted{\unsure{magasin|magasin,}} pré-atelier tel que

\[
X\ll Y\Longleftrightarrow X\leq Y\qquad X\mathrel{>\mkern-13mu<}Y\ \text{tjrs satisfaits.}
\]

de sorte que

\[
X\overset{\circ}{\ll}Y\Longleftrightarrow X=Y
\]

Alors tous les axiomes \inserted{Prat 1 - Prat 6} sont automatiquement satisfaits -- on peut donc partir d'un \deleted{\unsure{At|tel}} ens ordonné $(M,\leq)$ quelconque. Passant à $\unsure{\mathcal F|\mathscr F|F}$, on trouve \deleted{\unsure{c'est-à-dire|la totalité}}

\[
\begin{aligned}
&\mathcal F=P\!f(M,\leq)\qquad \text{ensemble de \uline{toutes}}\\
&\hspace{15em}\text{les parties fermées}
\end{aligned}
\]

% slanted
\margin{NB On trouve même un atelier \uline{séparé}, et \deleted{\unsure{disjoint}}\inserted{à multistrates près} \deleted{\illegible{1}} \unsure{ils|s'} sont (sauf erreur) \unsure{contrts|contr.} \unsure{//éditiles|//éditibles|//édités} etc\ldots}

% slanted
\margin{Pourtant, on peut ``faire de la topologie'' \unsure{cl$^{\mathrm{que}}$|classique|clque} sur une maquette simpliciale !}

dans $\mathcal F$, $F\ll G\Longleftrightarrow F\leq G$, et $F\mathrel{>\mkern-13mu<}G$ tjrs satisfait $F\parallel G\Longleftrightarrow F\cap G=\varnothing$. \deleted{\unsure{Axiome} C \unsure{est}} Les conditions $\mathrm{At}\,D$, $\mathrm{At}\,C(\unsure{\mathrm{spac}|\mathrm{spa}})$, $\mathrm{At}\,\mathrm{Fil}$, $\mathrm{At}\,\mathrm{supp}$ sont satisfaites\unsure{;|.} Pour ce \deleted{$\unsure{\mathrm{At}\,L\,1}$} qui est des conditions sur les lieux, $\mathrm{At}\,L\,2$ équivaut :

Tout \deleted{$\forall x\in M$} $X\in M$ est minoré par un élément minimal -- \unsure{c'est|n'est} pas \unsure{prédonnant|prédonné|précédent|prédécidé} ! L'axiome $\mathrm{At}\,L\,3$ tjrs satisfait (car deux lieux distincts sont disjoints, étant compatibles).

Donc le seul axiome pas satisfait est $\mathrm{At}\,L\,1$, divisibilité

\scanpage{55}
\mspage{54}

% slanted
\margin{\uline{25 Juin}}

Je voudrais revenir sur la notion d'\uline{atelier ensembliste}. Ce sera un atelier \unsure{soit} spacieux, soit polyédral, prenons le premier cas.

\uline{Proposition} Soit $\unsure{\mathcal A|\mathscr A|A}=(\unsure{\mathcal F|\mathscr F|F},\leq,\ll)$ un atelier, considérons le (co)morphisme canonique

\[
\unsure{\varphi|\phi}:\mathcal F\longrightarrow\mathrm{Fig}(\unsure{\mathcal L|\mathscr L})\qquad F\longmapsto\unsure{\mathrm{Malomb}|\mathrm{malomb}}(F)
\]

\inserted{Considérons les} Conditions : \deleted{\unsure{équivalentes}}

(i) L'hom. $\varphi$ est un \uline{plongement} d'ateliers (cf \unsure{CF|GF} \unsure{VI|II} p. 51) et l'image de $\mathcal F$ par $\varphi$ est un sous-atelier \uline{spacieux} (loc cit \inserted{p. 40} : si $F,G\in\varphi(\mathcal F)$ et $F\mathrel{>\mkern-13mu<}G$ \deleted{\unsure{entre|est}} dans $\mathrm{Fig}(\mathcal L)$, alors $F\mathrel{>\mkern-13mu<}G$ dans $\varphi(\mathcal F)$ i.e $F\cup G\in\deleted{\unsure{F}}\varphi(\mathcal F)$)

En d'autres termes :

% slanted
\margin{$\varphi$ plongement}

\[
\left\{\begin{aligned}
&\text{a) }\forall F\in\mathcal F,\quad F'\unsure{\mathrel{\vdash\mkern-6mu\dashv}|\longmapsto}\varphi(F')\ \text{induit une bijection}\\
&\hspace{6em}\mathrm{Ss\,fig}(F)\longrightarrow\mathrm{Ss\,fig}(\varphi(F))\\
&\text{b) }\varphi(F)\ll\varphi(G)\Rightarrow F\ll G\quad(\text{donc }\varphi\ \text{\unsure{injective|injection}})
\end{aligned}\right.
\]

% slanted
\margin{$\varphi$ plongement ``spacieux''}

\[
\text{c) Si }\varphi(F)\mathrel{>\mkern-13mu<}\varphi(G),\Rightarrow F\mathrel{>\mkern-13mu<}G
\]

(ii) $\mathcal F$ est \uline{spacieux} (au sens absolu) i.e. satisfait \deleted{\unsure{les}} l'un des quatre systèmes \unsure{équivalents|équivalent} de \unsure{conditions|condition} de la p. 45 \inserted{pour $\unsure{\Lambda|\wedge}=\mathcal L$} (p.ex. $\mathrm{At}\,L\,1$, $\mathrm{At}\,L\,2$, $\mathrm{At}\,L\,3$, $\mathrm{At}\,CL(\unsure{\mathrm{spac}|\mathrm{spa}})$), \uuline{et} pour $x,y\in\mathcal L$, $x\ne y$ \deleted{\unsure{on a}}\inserted{$\Rightarrow$} $x\deleted{\unsure{=}}\parallel y$. (donc $\mathrm{At}\,L\,3$ devient automatique)

Alors $(i)\Rightarrow(ii)$ et $(ii)\Rightarrow(i\,a)\deleted{\unsure{(b)|(c)}}$ (mais pas néc $(i\,b)$, \inserted{ni $(i,c)$,} ni l'injectivité de $\varphi$)

\uline{Dém} $(i)\Rightarrow(ii)$. \deleted{Comme} Commençons par $x,y\in\mathcal L$, $x\ne y\Rightarrow\deleted{\unsure{\varphi|y}}x\parallel y$ ou encore, \deleted{$(x,y)\in\mathcal L$}\inserted{$x,y\in\mathcal L$}$\Rightarrow x\mathrel{>\mkern-13mu<}y$. Comme $\varphi(x),\varphi(y)$ sont des \deleted{\unsure{familles}} figures ens. ponctuelles on a $\varphi(x)\mathrel{\deleted{\unsure{\times}}>\mkern-13mu<}\varphi(y)$ donc $x\mathrel{>\mkern-13mu<}y$ par c). $\mathrm{At}\,L\,1$ : Soit $X\in\unsure{M|\mathcal M}$, \deleted{soit $x$} comme $X\ne\varnothing$ on a $\varphi(\unsure{x|X})\ne\varnothing$, donc par définition de $\varphi$ \unsure{Cela|Celà} signifie $\mathrm{omb}(X)^{\circ}\ne\varnothing$, \deleted{\unsure{et|ce}} OK.

\scanpage{56}
\mspage{55}

$\mathrm{At}\,L\,2$ : Soient $X,Y\in\unsure{M|\mathcal M}$, avec $\forall x\in\mathrm{omb}(X),y\in\mathrm{omb}(Y)$ on a $x\ne y$ \inserted{($\Longleftrightarrow x\parallel y$)} i.e. $\mathrm{omb}(X)\cap\mathrm{omb}(Y)=\varnothing$. \deleted{Cela signifie que $\mathrm{omb}(X)\cap$ $\unsure{\varphi|\phi}(X)\parallel$} Comme $\mathrm{omb}(\unsure{x|X})=\lvert\varphi(x)\rvert$ et \unsure{item|itou|idem} pour $\mathrm{omb}(Y)$, on aura donc $\varphi(\unsure{X|x})\parallel\varphi(Y)$ a fortiori \deleted{$\mathrel{>\mkern-13mu<}$} $\varphi(X)\mathrel{>\mkern-13mu<}\varphi(Y)$ donc $X\mathrel{>\mkern-13mu<}Y$.

Il est évident d'autre part que $X\cap Y=\varnothing$ (à cause de $L\,1$ ; si $Z\leq X,Y$, $\exists x\in\unsure{\mathcal L|\mathscr L}\quad x\ll Z$, donc $x\in\mathrm{omb}(X)\cap\mathrm{omb}(Y)$, absurde).

$\mathrm{At}\,CL(\unsure{\mathrm{spac}|\mathrm{spa}})$ : Soient $X,Y\in M$, $L=\widetilde X\cap\widetilde Y$, supposons que $\forall x\in\mathrm{omb}(X)\setminus\mathrm{omb}(L),y\in\mathrm{omb}(Y)\setminus\mathrm{omb}(L)$ on ait $x\deleted{\unsure{=|\ne}}\parallel y$ \unsure{i.e.|i.e. $\in$.} $x\unsure{=|\ne}y$, en d'autres termes

\[
\mathrm{omb}(X)\cap\mathrm{omb}(Y)=\deleted{\varnothing}\mathrm{omb}(L)
\]

Prouvons $X\mathrel{>\mkern-13mu<}Y$. Or \deleted{on \unsure{aura|a}} \inserted{la} condition est

\[
\lvert\varphi(X)\rvert\cap\lvert\varphi(Y)\rvert=\lvert\varphi(L)\rvert
\]

où $\varphi(L)$ est une sous-figure ensembliste commune à $\varphi(X),\varphi(Y)$, d'où $\varphi(X)\mathrel{>\mkern-13mu<}\varphi(Y)$ donc $X\deleted{\unsure{\leq}}\mathrel{>\mkern-13mu<}Y$, OK.

$(ii)\Rightarrow(i\inserted{\,a,c})$ Il faut prouver a) \deleted{\unsure{b)}} c)

\uline{Lemme 1} Soit $\unsure{\mathcal A|\mathscr A|A}=(\unsure{\mathcal F|\mathscr F|F},\leq,\ll)$, \deleted{satisfaisant $L\,1$. Alors} \inserted{atelier, $\unsure{\Lambda|\wedge}$ partie dense de $M$}, \deleted{\unsure{Alors|telle}} Considérons \deleted{\unsure{l'hom.}} le comorphisme canonique

\[
\varphi_{\Lambda}:\mathcal F\longrightarrow\mathrm{Fig}(\Lambda)\qquad F\longmapsto\unsure{\mathrm{mulomb}|\mathrm{malomb}}_{\Lambda}(F)
\]

% slanted
\margin{La densité de $\Lambda$ sert à assurer que $\varphi_{\Lambda}$ transforme \unsure{les} strates en strates}

Alors $\varphi_{\Lambda}$ satisfait \unsure{c|a}, i.e. $\forall F$, on a

\[
\mathrm{Ss\,fig}(F)\xrightarrow{\sim}\unsure{\mathrm{ss\,fig}|\mathrm{sc\,fig}}(\varphi_{\Lambda}(F)),\quad F'\unsure{\mathrel{\vdash\mkern-6mu\dashv}|\longmapsto}\varphi_{\Lambda}(F')
\]

Il revient au $=$ de dire que l'application

\[
\widetilde F\longrightarrow\widetilde{\varphi_{\Lambda}(F)}\qquad \text{est un iso. d'ens}
\]

ordonnés

\scanpage{57}
\mspage{56}

\unsure{Ça|Çà} a été vu pour $\unsure{\Lambda|\wedge}=\unsure{M|\mathcal M}$. Dans le cas général, c'est un truc \unsure{ens|ens.}, ça marche pour applic. ens. quelc.

\[
i:\Lambda\longrightarrow M
\]

et figure ens. $\Phi\in\mathrm{Fig}(M)$, quand $\forall X\in\Phi$, on a $i^{\unsure{*|\times}}(X^{\circ})\ne\varnothing$.

Prenant $\Lambda=\unsure{\mathcal L|\mathscr L}$ dans \deleted{\unsure{notre}} le cas de la prop. \inserted{(supposant \unsure{(ii)|(i)})} on trouve a.

Prouvons \deleted{\unsure{b}}c). \deleted{\unsure{Le cas}} \inserted{Supposons $\unsure{\varphi|\phi}(F)\mathrel{>\mkern-13mu<}\varphi(G)$, prouvons $\deleted{\unsure{\varphi}}F\mathrel{>\mkern-13mu<}G$} Compte tenu de a) il s'ensuit que la plus grande sous-figure commune de $\varphi(F),\varphi(G)$ est $\varphi(L)$, où $L=F\cap G$ \deleted{$\unsure{\overset{\mathrm{déf}}{=}}\unsure{F\cap G|\mathrm{Inf}}$} ($\overset{\mathrm{déf}}{=}\mathrm{Inf}^{\leq}(F,G)$, donc $\widetilde L=\widetilde F\cap\widetilde G$). En effet, $\varphi(L)$ est une \inserted{telle} sous-figure commune, et si $\Phi$ en est \deleted{une autre}\inserted{la plus} grande telle sous-fig. commune de $\varphi(F),\varphi(G)$, on aura $\varphi(L)\unsure{\subset|\subseteq}\Phi$, \deleted{donc $\Phi$} et par a) on aura $\Phi=\varphi(L')=\varphi(L'')$, où $L'\leq F$, $L''\leq G$. \unsure{De|et} $L\leq L',L''$.

Prouvons que (compte tenu de $\varphi(L')=\varphi(L'')$) on a $L'=L''$ donc $L'=L''=L$. (\uline{NB} on ne sait pas si $\varphi$ est injectif, et $\varphi$ peut ne pas l'être\ldots)

Si c'était faux, soit $\unsure{X'|X''|X}$ \deleted{\unsure{autre}} $\in\deleted{\unsure{L}}\widetilde{L'}\setminus\unsure{\widetilde{L'}|\widetilde L,}$ il existe (because a)) un unique $\deleted{\unsure{X}}Y\unsure{^{\prime}|^{\prime\prime}}$ \inserted{$\in\widetilde{L''}\setminus\widetilde L$} telle que $\varphi(\deleted{\unsure{X'|Y}}\unsure{X|X'})=\varphi(\deleted{\unsure{Y'|X}}\unsure{Y|Y'})$. Alors la plus grande sous-figure commune de $\unsure{X'|X},\unsure{Y|Y'}$ est $X\cap L=Y\cap L$, distincte de $X,Y$. Est-ce bien impossible ? Si on trouve une telle situation, on aurait $\varphi(X)\mathrel{>\mkern-13mu<}\varphi(Y)$ puisque $\varphi(X)=\varphi(Y)$, mais $X\mathrel{>\mkern-13mu<}Y$ à cause de a) appliqué à \deleted{$\unsure{X\cup Y|X,Y}\unsure{\in F|F}$} $F=X\cup Y$ (si on avait $X\mathrel{>\mkern-13mu<}Y$).
